% =====================================================================
% PRIORITY REVISION PACK — ICCEREC 2026
% Paper: Adaptive Trust-Aware Access Control ... OpenDaylight-Based SDN
%
% IMPORTANT: These are REPLACEMENTS for existing text, not additions.
% Do NOT paste the sensitivity/robustness tables again (the manuscript
% already contains them; pasting would duplicate \label names).
%
% Contents:
%   1. Abstract (REPLACE)
%   2. Robustness subsection (REPLACE the closing paragraph)
%   3. Conclusion (REPLACE)
% =====================================================================


% ---------------------------------------------------------------------
% 1. ABSTRACT  (REPLACE the current abstract)
%    ~250 words; adds the measured enforcement performance (R3-4).
%    Update the latency/throughput numbers to your final canonical run.
% ---------------------------------------------------------------------
\begin{abstract}
Software-defined networking (SDN) enables centralized, programmable control, but
access decisions must adapt to changing user and device conditions instead of
trusting a network location. This study implements and evaluates a trust-aware
adaptive access control mechanism for Zero Trust security in an
OpenDaylight-based SDN. A trust score combines role/identity, contextual, and
behavioral factors through normalized weights and is mapped to three access
levels: FULL, LIMITED, and DENIED. Policy-specific OpenFlow rules are
provisioned through the OpenDaylight RESTCONF interface and enforced on a
Mininet--Open vSwitch data plane. In a four-switch ring, the three evaluated
scenarios yield baseline trust scores of 90, 69, and 38; FULL grants both
protected resources, LIMITED restricts access to the IoT service, and DENIED
blocks both, while logout returns the data plane to default-deny. Beyond
functional checks, the evaluation measures enforcement cost: the policy decision
point provisioned per-session rules in about 130~ms and revoked them in about
80~ms on average, and an authorized session sustained 71~Gbit/s through the
switch path, whereas the default-deny and post-logout states blocked all access.
A sensitivity analysis shows that the alternative weights 0.45/0.25/0.30
preserve all three classifications, that the threshold pairs 75/45 and 70/50
leave them unchanged, and that the mapping remains stable within a local
neighborhood of the chosen weights. The weights are experimental configurations
validated by sensitivity analysis rather than statistically optimized estimates.
Results are reported for a four-switch ring and establish functional correctness
together with measured enforcement performance.
\end{abstract}


% ---------------------------------------------------------------------
% 2. ROBUSTNESS SUBSECTION — REPLACE the closing paragraph.
%    Removes "sampling/grid resolution not specified" and states the
%    actual grid and counts, while keeping the honest scope.
% ---------------------------------------------------------------------
The preservation rates are computed from the closed-form trust model rather than
from data-plane measurements. Weights are enumerated on a grid of step $0.05$:
the full simplex comprises $231$ nonnegative, sum-to-one configurations, the
ordered domain ($w_R \ge w_C \ge w_B$) comprises $154$, and the local
neighborhood (each weight within $\pm 0.05$ of its baseline value) comprises
$25$. The rates are $27/231$ ($11.7\%$), $79/154$ ($51.3\%$), and $21/25$
($84.0\%$), respectively. Because the three evaluated scenarios lie close to the
$(70,40)$ boundaries, classification changes under remote weight assignments are
expected and confirm that the thresholds operate at the intended decision
boundary; near the chosen values the mapping is largely preserved. These rates
describe the assessed grid and are not statistical confidence levels or exact
continuous-volume fractions. The baseline weights therefore remain an
experimental configuration validated through sensitivity analysis, not a
statistically optimized solution.


% ---------------------------------------------------------------------
% 3. CONCLUSION — REPLACE the current conclusion.
%    Removes latency/controller load from future work (now measured) and
%    updates the future-work scope.
% ---------------------------------------------------------------------
The OpenDaylight-based prototype maps baseline trust scores of $90$, $69$, and
$38$ to FULL, LIMITED, and DENIED access, with matching resource permissions and
post-logout blocking. Beyond functional validation, the measured enforcement
cost is modest: per-session provisioning and revocation are on the order of
$100$~ms and $80$~ms, respectively, and an authorized session sustained
$71$~Gbit/s through the switch path, while default-deny and post-logout states
blocked all access. Sensitivity analysis identifies classification-preserving
weight and threshold alternatives together with boundary cases, and the reported
preservation rates support domain-dependent stability rather than robustness
throughout the weight space. The weights are experimental configurations
validated through sensitivity analysis, not statistically optimized estimates.
The evaluation does not establish enterprise-scale viability or attack
resilience; it validates the functional mapping from trust conditions to
resource-specific enforcement on a four-switch ring and reports controller-side
and data-plane performance. Future work will examine larger reference
topologies, continuous trust reassessment, and a learned behavioral factor that
augments the transparent trust aggregation.
